\section{AVO in brief}
\label{sec:avo}

This section summarizes the AVO paper \citep{chen2026avo} (arXiv v1, March 2026). We keep its
notation. Where the paper is silent, we say so, because those gaps are where an implementation must
make choices.

\subsection{Formulation}

Generic evolutionary search keeps a population $\Pop = \{(x_i, \f(x_i))\}$ of solutions and their
scores, and repeats
\begin{equation}
  \Pop_{t+1} = \Update\big(\Pt, (x_{t+1}, \f(x_{t+1}))\big), \qquad x_{t+1} = \Vary(\Pt).
  \label{eq:evo}
\end{equation}
Prior LLM methods split $\Vary$ into $\Generate \circ \Sample$ (\cref{eq:intro-generate}). The paper
notes that ``the LLM only participates in Generate: the sampling strategy, evaluation protocol,
population management, and the order of operations are all determined by the framework''
\citep[\S2.1]{chen2026avo}. AVO replaces the whole operator with one agent run
(\cref{eq:intro-agent}). The agent ``subsumes Sample, Generate, and evaluation into a single
autonomous loop''.

The objects are:
\begin{itemize}[leftmargin=*,itemsep=1pt]
  \item $\Pt = \{(x_1, \f(x_1)), \ldots, (x_t, \f(x_t))\}$, ``the full lineage of solutions and their
    scores''. In the experiments, each $x_i$ is a CUDA kernel with inline PTX.
  \item $\f(x_i) = (f_1(x_i), \ldots, f_n(x_i))$, a vector with one component per test
    configuration. A candidate that fails the correctness check gets $f_j(x_i) = 0$ for all $j$,
    ``regardless of throughput''.
  \item $\K$, a knowledge base: CUDA guides, the PTX ISA, Blackwell specifications, and existing
    kernels including the FlashAttention-4 source.
\end{itemize}

\subsection{One variation step}

A step that produces $x_{t+1}$ is an autonomous agent loop with ``numerous internal actions''. The
agent often reads several prior versions in $\Pt$, compares their profiles, and consults $\K$ before
it edits. It then calls $\f$. If the candidate fails or does not improve, the agent diagnoses and
revises. The paper calls this the \emph{edit--evaluate--diagnose} cycle. The rule for keeping a
result is stated once:
\begin{quote}\small
``we persist a new committed version only when it passes correctness checks and matches or improves
the benchmark score relative to the best committed version so far; unsuccessful intermediate
attempts remain part of the agent's internal search trajectory but are not added to the committed
lineage.'' \citep[\S3.2]{chen2026avo}
\end{quote}
Read literally, this is an elitist, monotone update: $\Pop_{t+1} = \Pt \cup \{(x', \f(x'))\}$ if
$x'$ is correct and scores at least as well as the best, and $\Pop_{t+1} = \Pt$ otherwise. Nothing is
pruned. The paper does not say how ``matches or improves'' applies to a \emph{vector}, or what
tolerance counts as a match. Its trajectory plots use the geometric mean across configurations,
which suggests a scalar reduction, but the text does not commit to one.

\subsection{Continuous evolution and the supervisor}

The paper runs AVO as a single lineage from a seed $x_0$, ``leaving population-level branching and
archive management to future extensions'' \citep[\S3.3]{chen2026avo}. Each committed version is a
git commit with its score. Two failure modes appear in long runs: the agent can \emph{stall} when it
exhausts a line of exploration, or it can enter ``unproductive cycles of edits that repeatedly fail
to improve scores''. A self-supervision mechanism detects both. When triggered, it ``reviews the
overall evolutionary trajectory and steers the search toward several candidate optimization
directions''. Figure~2 of the paper draws this as a separate \emph{AVO Supervisor Agent}. The paper
gives no thresholds, no trigger counts, and no mechanism for how the directions reach the main
agent.

The agent itself is ``an internally-developed general-purpose coding agent powered by frontier
LLMs'' with code editing, shell, file navigation and documentation retrieval. Its memory is its
conversation history. The models are not named, and no token or cost figures are given.

\subsection{Reported results}

\Cref{tab:avo-results} summarizes the results on B200 (forward prefill, head dimension 128, BF16,
sequence lengths 4K--32K at a fixed 32K total tokens). The seven-day run produced 40 versions and
explored ``over 500 candidate optimization directions''. Gains came in discrete jumps, for example
branchless accumulator rescaling ($+8.1\%$ non-causal at v20) and register rebalancing across warp
groups ($+2.1\%$ at v33). The paper uses these to argue that the agent performs ``genuine
hardware-level reasoning rather than superficial code transformations''. Adapting the final kernel
to grouped-query attention (GQA) took 30 more minutes of autonomous work.

\begin{table}[t]
\centering\small
\caption{Headline results of the AVO paper \citep[\S4]{chen2026avo}. Ranges are over the four
sequence lengths. TFLOPS values are the paper's bar labels.}
\label{tab:avo-results}
\begin{tabular}{@{}lccc@{}}
\toprule
Workload & Peak AVO TFLOPS & vs.\ cuDNN & vs.\ FlashAttention-4 \\
\midrule
MHA, causal      & 1637 & $+0.4\%$ to $+3.5\%$ & $+5.0\%$ to $+10.5\%$ \\
MHA, non-causal  & 1668 & $+1.8\%$ to $+2.4\%$ (16K, 32K) & within noise at 4K, 8K \\
GQA, causal      & 1647 & up to $+7.0\%$ & up to $+9.3\%$ \\
GQA, non-causal  & 1679 & up to $+6.0\%$ & up to $+4.5\%$ \\
\midrule
\multicolumn{4}{@{}l}{Run: 7 days, 40 committed versions, $>$500 directions explored, 1 lineage.}\\
\bottomrule
\end{tabular}
\end{table}

\subsection{What the paper leaves open}
\label{sec:avo-open}

For an implementer, the important gaps are:
\begin{enumerate}[leftmargin=*,itemsep=1pt]
  \item \textbf{No code or harness.} No repository, prompts, or agent configuration.
  \item \textbf{An underspecified acceptance rule.} ``Matches or improves'' on a vector, with no
    noise band, although the authors measured each benchmark 10 times for mean and deviation.
  \item \textbf{An unspecified supervisor.} No detection rule, threshold, or delivery mechanism.
  \item \textbf{Trajectory is lost.} Failed attempts live only in the agent's context, so a new
    session cannot learn from them.
  \item \textbf{No population.} Branching and archives are future work.
  \item \textbf{No ablations of the harness.} Neither $\K$ nor the supervisor is ablated, and AVO is
    not compared with a $\Generate \circ \Sample$ baseline on the same budget.
\end{enumerate}
avocode gives a concrete answer to items 1--4 and a partial answer to item 5. It does not address
item 6. \Cref{sec:discussion} returns to each.
